%%%%%%%%%%%%%%%%%%%%%%%%%%%%%%%%%%%%%%%%%
% Presentasi Teori: Metodologi Quantum Bayesian Game Theory
% Berdasarkan Adv_Gemini_plan.md
%%%%%%%%%%%%%%%%%%%%%%%%%%%%%%%%%%%%%%%%%

\documentclass[aspectratio=169]{beamer}
\usepackage[utf8]{inputenc}
\usepackage{amsmath,amssymb,physics}
\usepackage{graphicx}
\usepackage{tikz}
\usepackage{booktabs}
\usepackage{hyperref}

% Theme
\usetheme{Madrid}
\usecolortheme{seahorse}

% Judul
\title{Metodologi Quantum Bayesian Game Theory untuk Analisis Pasar}
\subtitle{Dari Econophysics ke Quantum Entanglement}
\author{Presentasi Tugas Akhir}
\date{\today}

\begin{document}

% Slide Judul
\begin{frame}
    \titlepage
\end{frame}

% Outline
\begin{frame}{Outline}
    \tableofcontents
\end{frame}

%%%%%%%%%%%%%%%%%%%%%%%%%%%%%%%%%%%%%%%%%
\section{Pendahuluan}
%%%%%%%%%%%%%%%%%%%%%%%%%%%%%%%%%%%%%%%%%

\begin{frame}{Pendahuluan}
    \begin{itemize}
        \item Metodologi yang menjembatani dinamika pasar (\textbf{Econophysics}) dengan interaksi strategis (\textbf{Game Theory})
        \item Menggunakan formalisme mekanika kuantum
        \item Menangkap ``keterikatan'' antar aset melalui probabilitas keadaan bersama
        \item Tidak sekadar korelasi linear, melainkan representasi \textbf{quantum state}
    \end{itemize}
\end{frame}

%%%%%%%%%%%%%%%%%%%%%%%%%%%%%%%%%%%%%%%%%
\section{Konstruksi Matriks Payoff}
%%%%%%%%%%%%%%%%%%%%%%%%%%%%%%%%%%%%%%%%%

\begin{frame}{Skema Leader-Follower (Stackelberg)}
    \begin{itemize}
        \item Pergerakan harga diperlakukan sebagai strategi
        \item Notasi bra-ket $\ket{\cdot}$ untuk merepresentasikan \textit{state} diskrit (naik/turun)
    \end{itemize}
    
    \vspace{0.5cm}
    
    \begin{table}
        \centering
        \begin{tabular}{c|c|c}
            \toprule
            \textbf{Leader / Follower} & $\ket{0}_F$ (Up) & $\ket{1}_F$ (Down) \\
            \midrule
            $\ket{0}_L$ (Up) & $(a, b)$ & $(c, d)$ \\
            $\ket{1}_L$ (Down) & $(e, f)$ & $(g, h)$ \\
            \bottomrule
        \end{tabular}
    \end{table}
\end{frame}

\begin{frame}{Rumus Nilai Payoff}
    Nilai payoff dihitung menggunakan rata-rata return dengan Laplace smoothing:
    
    \begin{align}
        a &= \frac{\sum L^{00}}{n_{00} + 1}, \quad b = \frac{\sum F^{00}}{n_{00} + 1} \\[10pt]
        c &= \frac{\sum L^{01}}{n_{01} + 1}, \quad d = \frac{\sum F^{01}}{n_{01} + 1} \\[10pt]
        e &= \frac{\sum L^{10}}{n_{10} + 1}, \quad f = \frac{\sum F^{10}}{n_{10} + 1} \\[10pt]
        g &= \frac{\sum L^{11}}{n_{11} + 1}, \quad h = \frac{\sum F^{11}}{n_{11} + 1}
    \end{align}
    
    \vspace{0.3cm}
    Dimana $n_{ij}$ adalah jumlah kejadian state $(i,j)$
\end{frame}

%%%%%%%%%%%%%%%%%%%%%%%%%%%%%%%%%%%%%%%%%
\section{Pemodelan Game Theory}
%%%%%%%%%%%%%%%%%%%%%%%%%%%%%%%%%%%%%%%%%

\begin{frame}{Klasifikasi Tipe Permainan}
    Berdasarkan nilai payoff, hubungan antar aset dapat diklasifikasikan:
    
    \begin{itemize}
        \item \textbf{Coordination Game}: Payoff tertinggi pada $(0,0)$ dan $(1,1)$
            \begin{itemize}
                \item Aset bergerak searah (korelasi positif kuat)
            \end{itemize}
        \item \textbf{Anti-Coordination}: Payoff tertinggi pada $(0,1)$ atau $(1,0)$
            \begin{itemize}
                \item Hubungan substitusi atau \textit{hedging}
            \end{itemize}
        \item \textbf{Prisoner's Dilemma}: Ada insentif ``mengkhianati'' arah tren
    \end{itemize}
    
    \vspace{0.5cm}
    \textbf{Nash Equilibrium}: Jika $a > e$ dan $b > d$, maka $\ket{00}$ adalah equilibrium
\end{frame}

\begin{frame}{Bayesian Game}
    \begin{itemize}
        \item Matriks dapat menjadi \textbf{Bayesian Game}
        \item \textbf{Type}: Kondisi internal pasar (bullish/bearish tersembunyi)
        \item \textbf{Beliefs}: Probabilitas pergerakan lawan
    \end{itemize}
    
    \vspace{0.5cm}
    
    Keadaan $\ket{\psi}$ merepresentasikan \textit{Joint Probability Distribution} yang mendasari \textbf{Bayesian Nash Equilibrium}
\end{frame}

%%%%%%%%%%%%%%%%%%%%%%%%%%%%%%%%%%%%%%%%%
\section{Quantum Entanglement}
%%%%%%%%%%%%%%%%%%%%%%%%%%%%%%%%%%%%%%%%%

\begin{frame}{Konstruksi Quantum State}
    Sistem dua aset direpresentasikan sebagai fungsi gelombang:
    
    \begin{equation}
        \ket{\psi} = a_{00} \ket{00} + a_{01} \ket{01} + a_{10} \ket{10} + a_{11} \ket{11}
    \end{equation}
    
    \vspace{0.5cm}
    
    \begin{itemize}
        \item $|a_{ij}|^2 = P(\ket{ij})$ adalah probabilitas gabungan
        \item Koefisien dihitung dari: $a_{ij} = \sqrt{P_{ij}}$
    \end{itemize}
\end{frame}

\begin{frame}{Reduced Density Matrix \& Entanglement}
    Untuk mengukur derajat ketergantungan (entanglement):
    
    \begin{enumerate}
        \item Bentuk matriks densitas: $\rho = \ket{\psi}\bra{\psi}$
        \item Partial trace: $\rho_L = \text{Tr}_F(\rho)$
        \item Von Neumann Entropy:
        \begin{equation}
            S = -\text{Tr}(\rho_L \ln \rho_L)
        \end{equation}
    \end{enumerate}
    
    \vspace{0.5cm}
    
    \textbf{Interpretasi}: Jika $S > 0$, kedua aset \textit{ter-entangle} secara informasi
\end{frame}

\begin{frame}{Quantum Mutual Information}
    \begin{equation}
        I(L:F) = S(\rho_L) + S(\rho_F) - S(\rho_{LF})
    \end{equation}
    
    \vspace{0.5cm}
    
    \begin{itemize}
        \item Semakin tinggi nilai $I$, semakin kuat keterikatan
        \item Menunjukkan ``kluster kuantum'' di pasar keuangan
        \item Aset bergerak sebagai satu kesatuan fungsi gelombang
    \end{itemize}
\end{frame}

%%%%%%%%%%%%%%%%%%%%%%%%%%%%%%%%%%%%%%%%%
\section{Network Theory}
%%%%%%%%%%%%%%%%%%%%%%%%%%%%%%%%%%%%%%%%%

\begin{frame}{Quantum Financial Network}
    Untuk lebih dari 2 aset:
    
    \begin{itemize}
        \item \textbf{Nodes}: Setiap aset
        \item \textbf{Edges}: Nilai Entanglement Entropy atau Quantum Mutual Information
    \end{itemize}
    
    \vspace{0.5cm}
    
    Metode ini memungkinkan:
    \begin{itemize}
        \item Melihat siapa yang berkorelasi dengan siapa
        \item Menemukan ``penggerak utama'' (Leader) menggunakan \textbf{Eigenvector Centrality}
    \end{itemize}
\end{frame}

%%%%%%%%%%%%%%%%%%%%%%%%%%%%%%%%%%%%%%%%%
\section{Konstruksi Hamiltonian}
%%%%%%%%%%%%%%%%%%%%%%%%%%%%%%%%%%%%%%%%%

\begin{frame}{Model Ising Hamiltonian}
    Hamiltonian merepresentasikan \textbf{total energi} atau \textbf{fungsi biaya} sistem:
    
    \begin{equation}
        H = -\sum h_i \sigma_i^z - \sum J_{ij} \sigma_i^z \sigma_j^z
    \end{equation}
    
    \vspace{0.3cm}
    
    Komponen:
    \begin{itemize}
        \item \textbf{Field Term}: $H_{bias} = -\sum_{i} h_i \sigma_i^z$
            \begin{itemize}
                \item $h_i$ dari selisih payoff naik vs turun
            \end{itemize}
        \item \textbf{Interaction Term}: $H_{int} = -J_{LF} (\sigma_L^z \sigma_F^z)$
            \begin{itemize}
                \item $J_{LF} > 0$: Coordination Game (Feromagnetik)
                \item $J_{LF} < 0$: Anti-coordination
            \end{itemize}
    \end{itemize}
\end{frame}

\begin{frame}{Pemetaan Parameter}
    Menentukan bias ($h$) dari payoff:
    \begin{equation}
        h_L \approx \frac{(a+c) - (e+g)}{2}
    \end{equation}
    
    \vspace{0.3cm}
    
    Menentukan interaksi ($J$) dari Quantum Mutual Information:
    \begin{equation}
        J_{LF} = \text{sign(korelasi)} \times I(L:F)
    \end{equation}
\end{frame}

\begin{frame}{Aplikasi Hamiltonian}
    \begin{enumerate}
        \item \textbf{Evolusi Waktu} (Schrödinger Equation):
        \begin{equation}
            \ket{\psi(t+dt)} = e^{-iHdt} \ket{\psi(t)}
        \end{equation}
        
        \item \textbf{Ground State}: Prediksi posisi pasar paling stabil
            \begin{itemize}
                \item Analog dengan Nash Equilibrium dalam ruang Hilbert
            \end{itemize}
        
        \item \textbf{Phase Transition}: Menghitung ``suhu kritis'' pasar saat krisis
    \end{enumerate}
\end{frame}

%%%%%%%%%%%%%%%%%%%%%%%%%%%%%%%%%%%%%%%%%
\section{Kesimpulan}
%%%%%%%%%%%%%%%%%%%%%%%%%%%%%%%%%%%%%%%%%

\begin{frame}{Kesimpulan}
    Metodologi lengkap dari:
    \begin{enumerate}
        \item Data mentah pergerakan aset
        \item Matriks payoff Game Theory
        \item Perhitungan Quantum Entanglement
        \item Pemetaan Network Theory
        \item Konstruksi Hamiltonian untuk prediksi
    \end{enumerate}
    
    \vspace{0.5cm}
    
    Transformasi dari analisis data historis menjadi model \textbf{mekanika kuantum stokastik} yang utuh
\end{frame}

\end{document}
